% =============================================================================
% NeurIPS 2026 submission - Main paper
%
% TODO(style-file): This draft currently loads `neurips_2025.sty`, which is the
% latest mirror available on public GitHub at the time of scaffolding. The
% official `neurips_2026.sty` is gated behind authenticated NeurIPS download
% pages. Before submission, drop the official `neurips_2026.sty` (and the
% matching `neurips_2026.tex` skeleton) into this directory and change the
% \usepackage line below from `neurips_2025` to `neurips_2026`. The 2026 file
% is structurally identical to 2025 in every recent year.
% =============================================================================

\documentclass{article}

% if you need to pass options to natbib, use, e.g.:
\PassOptionsToPackage{numbers, compress}{natbib}

% Submission (anonymous) - default option
\usepackage{neurips_2025}

% Preprint (e.g. for arXiv) - swap to this when going public
% \usepackage[preprint]{neurips_2025}

% Camera-ready - swap to this only after acceptance
% \usepackage[final]{neurips_2025}

\usepackage[utf8]{inputenc}
\usepackage[T1]{fontenc}
\usepackage{hyperref}
\usepackage{url}
\usepackage{booktabs}
\usepackage{amsfonts}
\usepackage{amsmath}
\usepackage{amssymb}
\usepackage{nicefrac}
\usepackage{microtype}
\usepackage{xcolor}
\usepackage{graphicx}
\usepackage{multirow}
\usepackage{subcaption}

% =============================================================================
% TODO(title): Final title TBD. Working title below. Aim < 12 words.
% Candidates:
%   - "Predicting Anisotropic Task-Vector Coefficients from Text and a Few Shots"
%   - "Multi-Modal Hypernetworks for Zero- and Few-Shot Task Vector Composition"
%   - "From Description to Coefficients: Hypernetwork-Predicted aTLAS for CLIP"
% =============================================================================
\title{Predicting Anisotropic Task-Vector Coefficients from Text and a Few Shots}

% =============================================================================
% TODO(authors): Fill in real author block before submission.
% Reminder: NeurIPS submission is double-blind; keep this anonymized in the
% submission build. The default option (no [final]/[preprint]) anonymizes for us.
% =============================================================================
\author{%
  Anonymous Authors \\
  Anonymous Affiliation \\
  \texttt{anonymous@example.com}
}

\begin{document}

\maketitle

% =============================================================================
% ABSTRACT
% TODO: Write last. ~150-200 words. Structure:
%   1. Problem: adapting CLIP-class models to new tasks needs labeled data; aTLAS
%      reduces it but still requires per-task coefficient learning.
%   2. Approach: meta-train a hypernetwork that predicts aTLAS coefficients from
%      a text description (zero-shot) or text + K support images (few-shot).
%   3. Method highlights: dual-branch fusion, variable-shot meta-training,
%      graceful text-only fallback.
%   4. Results: headline numbers on held-out datasets at K in {0,1,2,4,8,16}.
%   5. Takeaway: matches/beats per-task aTLAS at a fraction of the compute.
% =============================================================================
\begin{abstract}
% TODO(abstract): Write final pass after results are locked.
\end{abstract}

% =============================================================================
\section{Introduction}
% =============================================================================
% TODO(intro): ~1 page. Hit these beats in order:
%   - Setup: CLIP zero-shot is convenient but trails fine-tuning; per-task
%     fine-tuning is expensive.
%   - Task vectors / task arithmetic give a cheap composition primitive.
%   - aTLAS adds blockwise anisotropic scaling; learns coefficients per task.
%   - Limitation: aTLAS still needs labeled adaptation data per task.
%   - Our contribution: predict the coefficients directly, conditioned on a
%     text description (and optionally K support images) - turns aTLAS into a
%     genuinely zero-/few-shot method.
%   - Bullet list of contributions:
%       (i) text-to-coef hypernetwork formulation;
%       (ii) multi-modal extension with text+image fusion;
%       (iii) meta-training recipe with variable-K shots;
%       (iv) empirical: matches per-task aTLAS, beats CLIP zero-shot, beats
%            adapter baselines at low K.
%   - Paper roadmap.

% =============================================================================
\section{Related Work}
% =============================================================================

\subsection{Task Vectors and Task Arithmetic}
% TODO: Ilharco et al. 2023 (task arithmetic), Ortiz-Jimenez et al. (tangent
% task arithmetic), Yadav et al. (TIES-merging), task vector negation, etc.
% Frame: task vectors are scalar-weighted; we generalize to learned
% block-anisotropic weights and predict them.

\subsection{Anisotropic Scaling and aTLAS}
% TODO: aTLAS (the prior paper this builds on). Briefly recap the contribution
% and clarify what is and is not borrowed. Block coefficients, partition
% variant (aTLAS x K), linearized variant.

\subsection{Hypernetworks for Adaptation}
% TODO: Ha et al. 2017, HyperNets for few-shot, HyperLoRA, etc. Distinguish:
% prior hypernet work predicts weight matrices/LoRA; we predict task-vector
% coefficients - far smaller output space, much cheaper to train.

\subsection{Few-Shot and Zero-Shot Adaptation of CLIP}
% TODO: CoOp, CoCoOp, Tip-Adapter (and Tip-F), LP++, WiSE-FT. We compare
% directly to these. Note that these methods either learn prompts (CoOp) or
% adapters; we learn a different primitive entirely.

\subsection{Meta-Learning}
% TODO: MAML and descendants briefly; ProtoNet for K-shot conditioning. Our
% recipe is closer to amortized inference than gradient-based meta-learning.

% =============================================================================
\section{Background}
% =============================================================================

\subsection{CLIP and Task Vectors}
% TODO: Define notation. Pretrained encoder f_theta_0; fine-tuned f_theta_t for
% task t; task vector tau_t = theta_t - theta_0.

\subsection{Anisotropic Task-Level Scaling (aTLAS)}
% TODO: Block decomposition theta = (theta^{(1)}, ..., theta^{(B)}). aTLAS
% composes K task vectors with per-block coefficients
% \alpha \in R^{T x B}: theta_new = theta_0 + sum_t sum_b alpha_{t,b} tau_t^{(b)}.
% Standard aTLAS learns alpha per target task via gradient descent on labeled
% data; our setting amortizes this.

% =============================================================================
\section{Method}
% =============================================================================

\subsection{Problem Formulation}
% TODO: Given a target task described by (i) a text description set D and
% (ii) optionally K labeled support images per class, predict alpha so that
% the composed encoder performs well on held-out test data of the task.
% Key constraint: the prediction function is task-agnostic (single hypernet
% used for all tasks at test time, no task-specific gradient updates).

\subsection{Text-to-Coefficient Hypernetwork}
% TODO: Architecture: CLIP text encoder (frozen) -> aggregate over descriptions
% (mean) -> MLP -> alpha. Discuss aggregation: per-class vs. dataset-level.

\subsection{Multi-Modal Hypernetwork}
% TODO: Add a parallel image branch: OpenCLIP visual encoder (frozen) -> shot
% pooling (mean or attention) over K shots per class -> projection -> fuse with
% text projection. Three fusion modes: concat / add / cross-attention.
% Graceful fallback: when K=0, route text projection through a learned
% text_only_proj layer to keep dimensions consistent.

\subsection{Meta-Training Procedure}
% TODO: Episode = (task t, text D_t, K-shot support S_t, query batch Q_t).
% Forward pass produces alpha_hat = h_phi(D_t, S_t); compose model;
% cross-entropy loss on Q_t; gradient flows back through alpha_hat into phi.
% Variable-K trick: at each episode randomly sample K in {0, 1, 2, 4, 8, 16}
% (or uniform in [1, K_max]) so a single hypernetwork covers the full range.

\subsection{Inference}
% TODO: Forward h_phi once at test time, compose model, evaluate. No
% per-task gradient steps. Compare costs to per-task aTLAS / Tip-F.

% =============================================================================
\section{Experiments}
% =============================================================================

\subsection{Experimental Setup}
% TODO:
%   - Backbones: ViT-B/32, ViT-B/16, ViT-L/14 (HuggingFace CLIP).
%   - Meta-train tasks: CIFAR10, EuroSAT, DTD, GTSRB, SVHN, Food101.
%   - Held-out (eval) tasks: Caltech101, Flowers102, Cars, RESISC45, SUN397.
%   - K in {0, 1, 2, 4, 8, 16}; 3 seeds; report mean +/- std.
%   - Baselines: CLIP zero-shot; per-task aTLAS; Tip-Adapter; Tip-Adapter-F;
%     LP++; CoOp; WiSE-FT.
%   - Hyperparameters: see appendix.

\subsection{Main Results}
% TODO: Headline table - rows = methods, columns = K levels, cells = mean
% accuracy across held-out datasets +/- std over 3 seeds. Bold the best.
% Companion figure: line plot of accuracy vs. K, one curve per method.
% Per-dataset breakdown moves to appendix.

\subsection{Multi-Modal vs. Text-Only}
% TODO: Direct comparison table at K in {0, 1, 2, 4, 8, 16}. Expect
% multi-modal > text-only by 5-15% at low K, gap shrinks as K grows.

\subsection{Ablations}

\subsubsection{Fusion Strategy}
% TODO: concat vs. add vs. cross-attention. Single table; one number per
% (fusion, K) combination.

\subsubsection{Shot Pooling}
% TODO: mean vs. attention pooling.

\subsubsection{Hypernetwork Capacity}
% TODO: small / medium / large architectures from CLAUDE.md.

\subsubsection{Text Source}
% TODO: manual vs. CLIP-style templates vs. GPT-4o-generated vs.
% Claude-generated descriptions.

\subsubsection{Variable-K Meta-Training}
% TODO: training on a single K vs. variable K; show variable-K is robust
% across all eval K.

\subsection{Cross-Architecture Generalization}
% TODO: train hypernetwork on ViT-B/32, evaluate on ViT-B/16 and ViT-L/14.
% Even partial transfer is a strong selling point. May go to appendix if time
% is short.

\subsection{Compute Analysis}
% TODO: wall-clock + FLOPs comparison: per-task aTLAS vs. our hypernetwork at
% inference. Argue strong efficiency advantage.

% =============================================================================
\section{Discussion and Limitations}
% =============================================================================
% TODO:
%   - When does multi-modal help most? (low K, visually distinctive tasks.)
%   - When does it not help? (text-only sufficient for semantically clean
%     tasks; failures when text descriptions are misleading.)
%   - Limitations: tied to the meta-train dataset distribution; predicts only
%     scalar coefficients (not full LoRA); domain shift at test time.
%   - Broader impact: standard CLIP-class concerns (bias, dual use). Brief.

% =============================================================================
\section{Conclusion}
% =============================================================================
% TODO: Two-paragraph wrap. Restate contribution; future work (full LoRA
% prediction, larger meta-train pools, cross-modal applications).

% =============================================================================
% References
% =============================================================================
% TODO(refs): Maintain references.bib in this directory. Cite at minimum:
%   - Radford et al. 2021 (CLIP)
%   - Wortsman et al. 2022 (WiSE-FT)
%   - Ilharco et al. 2023 (task arithmetic)
%   - Ortiz-Jimenez et al. 2023 (tangent task arithmetic)
%   - aTLAS (NeurIPS 2024 - this paper's predecessor)
%   - Ha et al. 2017 (hypernetworks)
%   - Zhou et al. (CoOp / CoCoOp)
%   - Zhang et al. (Tip-Adapter)
%   - Huang et al. (LP++)
%   - Yadav et al. (TIES-merging)
\bibliographystyle{plainnat}
\bibliography{references}

% =============================================================================
\appendix
% =============================================================================

\section{Implementation Details}
% TODO: full hyperparameters (learning rates, schedules, batch sizes, optimizer,
% mixed precision settings). Hardware. Random seed protocol.

\section{Dataset Details}
% TODO: per-dataset summary; train/val/test sizes; meta-train vs. eval split
% rationale. Mostly a table.

\section{Per-Dataset Results}
% TODO: full breakdown of main table by dataset, all methods, all K.

\section{Additional Ablations}
% TODO: anything that didn't fit in the main paper:
%   - sensitivity to number of text descriptions per class
%   - effect of freezing vs. unfreezing the text/image encoders
%   - synthetic-data variant results (if included)

\section{Failure Cases}
% TODO: 1-2 qualitative examples where multi-modal underperforms text-only,
% with brief analysis.

% =============================================================================
% NeurIPS Paper Checklist (mandatory; absence is grounds for desk reject)
% =============================================================================
% TODO(checklist): The official 2026 .sty file ships with the checklist macros.
% Once we swap the sty in (see TODO at top of file), copy the official
% checklist template here and answer every item. Do this LAST, but do NOT
% skip it.

\end{document}
